% =========================================================================
% CHAPTER 3: MÔ HÌNH NGHIÊN CỨU VÀ THIẾT KẾ ĐỀ XUẤT
% =========================================================================

\section{Mô hình nghiên cứu và Thiết kế đề xuất}

% --- SLIDE 3.1: Yêu cầu thiết kế ---
\begin{frame}{3.1 Yêu cầu đối với mô hình đề xuất}
    \begin{columns}[T]
        \column{0.48\textwidth}
        \textbf{Yêu cầu chức năng:}
        \begin{itemize}
            \item Kết nối người dùng và dịch vụ theo từng khu vực chức năng.
            \item Phân tách lưu lượng giữa các khoa, phòng ban và vùng máy chủ.
            \item Kết nối Campus chính với các cơ sở phụ.
            \item Hỗ trợ quản lý tập trung tại LAN và WAN.
        \end{itemize}

        \column{0.48\textwidth}
        \textbf{Yêu cầu kiến trúc:}
        \begin{itemize}
            \item Phân tầng rõ ràng, hạn chế phụ thuộc giữa các khối.
            \item Có đường truyền và gateway dự phòng ở vị trí quan trọng.
            \item Tách vùng người dùng, quản trị, máy chủ và DMZ.
            \item Cho phép biểu diễn đầy đủ trên môi trường mô phỏng.
        \end{itemize}
    \end{columns}
\end{frame}

% --- SLIDE 3.2: Kiến trúc tổng thể ---
\begin{frame}{3.2 Kiến trúc tổng thể đề xuất}
    \begin{columns}[T]
        \column{0.31\textwidth}
        \begin{block}{Campus chính}
            \centering
            Core--Distribution--Access\\
            SDN Controller\\
            Server Farm và DMZ
        \end{block}

        \column{0.31\textwidth}
        \begin{block}{Kết nối liên cơ sở}
            \centering
            WAN Edge\\
            Internet và MPLS\\
            SD-WAN Overlay
        \end{block}

        \column{0.31\textwidth}
        \begin{block}{Các cơ sở phụ}
            \centering
            Site 200 -- Cần Thơ\\
            Site 300 -- Đà Nẵng\\
            Site 400 -- Nha Trang
        \end{block}
    \end{columns}
    \vspace{0.2cm}
    \begin{block}{Nguyên tắc tích hợp}
        SDN điều khiển miền chuyển mạch trong Campus; SD-WAN đảm nhiệm kết nối giữa các site. Hai miền gặp nhau tại lớp Core/WAN Edge.
    \end{block}
\end{frame}

% --- SLIDE 3.3: Mô hình Campus chính ---
\begin{frame}{3.3 Mô hình Campus chính}
    \begin{columns}[T]
        \column{0.48\textwidth}
        \textbf{Kiến trúc phân tầng:}
        \begin{itemize}
            \item \textbf{Core:} kết nối các vùng chính, định tuyến và cung cấp gateway.
            \item \textbf{Distribution:} tổng hợp các nhánh Access và áp dụng chính sách.
            \item \textbf{Access:} kết nối thiết bị người dùng theo VLAN.
        \end{itemize}

        \column{0.48\textwidth}
        \textbf{Vị trí của SDN:}
        \begin{itemize}
            \item Controller cung cấp lớp điều khiển logic tập trung.
            \item Các switch Distribution/Access đảm nhiệm chuyển tiếp dữ liệu.
            \item Kênh điều khiển được tách khỏi lưu lượng người dùng.
        \end{itemize}
    \end{columns}
    \vspace{0.15cm}
    \begin{block}{Các vùng dịch vụ}
        Server Farm cung cấp dịch vụ nội bộ; DMZ chứa dịch vụ cần công bố; Firewall kiểm soát lưu lượng giữa các vùng.
    \end{block}
\end{frame}

% --- SLIDE 3.4: Phân vùng logic ---
\begin{frame}{3.4 Phân vùng logic và quy hoạch địa chỉ}
    \begin{table}
        \centering
        \small
        \renewcommand{\arraystretch}{1.3}
        \begin{tabular}{|p{0.22\textwidth}|p{0.3\textwidth}|p{0.38\textwidth}|}
            \hline
            \rowcolor{primaryDark!15}
            \textbf{Nhóm mạng} & \textbf{Đối tượng} & \textbf{Nguyên tắc thiết kế} \\
            \hline
            Người dùng & Khoa, phòng ban, phòng học & Một VLAN/subnet riêng cho mỗi nhóm \\
            \hline
            Dịch vụ nội bộ & DHCP, Syslog và dịch vụ dùng chung & Đặt trong Server Farm, địa chỉ ổn định \\
            \hline
            Dịch vụ công bố & Web, Mail & Tách riêng trong DMZ và đi qua Firewall \\
            \hline
            Quản trị & Controller và giao diện quản lý & Tách khỏi lưu lượng người dùng \\
            \hline
            Liên kết WAN & Các đường point-to-point & Dùng dải riêng cho underlay và định danh site \\
            \hline
        \end{tabular}
    \end{table}
\end{frame}

% --- SLIDE 3.5: Mô hình liên kết đa cơ sở ---
\begin{frame}{3.5 Mô hình kết nối đa cơ sở}
    \begin{columns}[T]
        \column{0.48\textwidth}
        \textbf{Underlay:}
        \begin{itemize}
            \item Internet và MPLS cung cấp khả năng truyền dẫn vật lý.
            \item Mỗi site kết nối vào WAN thông qua thiết bị biên.
            \item Định tuyến underlay bảo đảm khả năng tiếp cận giữa các WAN Edge.
        \end{itemize}

        \column{0.48\textwidth}
        \textbf{Overlay:}
        \begin{itemize}
            \item Các site trao đổi route logic thông qua lớp điều khiển SD-WAN.
            \item Dữ liệu liên cơ sở được vận chuyển qua tunnel bảo mật.
            \item Chính sách WAN được quản lý theo site và loại lưu lượng.
        \end{itemize}
    \end{columns}
    \vspace{0.15cm}
    \begin{block}{Phạm vi của mô hình}
        Campus chính và ba cơ sở phụ được biểu diễn trong cùng một topology để nghiên cứu quan hệ giữa mạng nội bộ và mạng WAN.
    \end{block}
\end{frame}
